\section{Objetivo actualizado y medible}

El objetivo general responde a la pregunta de la sección~2.1:

\begin{quote}
\emph{Evaluar en qué medida las circunstancias registradas al levantar el parte permiten clasificar retrospectivamente cuáles de los 17.236 siniestros registrados por la ANT entre el 1 de enero y el 30 de septiembre de 2026 tuvieron al menos un fallecido en el lugar de los hechos, comparando modelos de aprendizaje automático clásico con referencias sencillas en casos posteriores en el tiempo, para conocer qué información del parte se asocia con la gravedad y complementar las estadísticas existentes.}
\end{quote}

\begin{description}[leftmargin=1.6em,style=nextline]
\item[OE1. Caracterizar la variabilidad de la gravedad.] Caracterizar cómo varía la proporción de siniestros graves entre grupos definidos con el conjunto~A, informando para cada grupo el número de casos y una medida de incertidumbre, para identificar qué circunstancias se asocian con mayor gravedad. Se comprueba con el criterio CA-03.

\item[OE2. Evaluar la clasificación.] Evaluar la clasificación retrospectiva de los siniestros graves con el conjunto~A, comparando al menos cinco tipos de modelos de aprendizaje automático clásico con dos referencias sencillas ---la proporción global y una regla basada solo en el tipo de siniestro--- en casos posteriores en el tiempo, para determinar si superan a reglas simples y cuánto añade la causa probable del conjunto~B. Se comprueba con los criterios CA-02 a CA-07.

\item[OE3. Comunicar el resultado.] Comunicar el resultado del modelo seleccionado mediante una consulta académica que devuelva la probabilidad estimada, las circunstancias asociadas y una advertencia de uso retrospectivo, para que usuarios no especialistas lo interpreten sin usarlo como alerta. Se comprueba con el criterio CA-09.
\end{description}
